% UBT-AI-PROVENANCE-BEGIN
% schema: ubt-ai-provenance/v1
% tier: C_working
% ai_assistance: disclosed
% human_review: risk-based
% editorial_responsibility: Ing. David Jaroš
% policy: ../../../AI_PROVENANCE.md
% notice: Working material; exhaustive human review is not claimed.
% UBT-AI-PROVENANCE-END

\chapter{Kovariantní derivace, konexe a čistá geometrie UBT}
\label{ch:covariant-tetrad}

Tato kapitola vysvětluje geometrické jádro UBT ve formě, která je nejbližší
původní intuici teorie:
\[
 \Theta\quad\longrightarrow\quad E_\mu:=D_\mu\Theta
 \quad\longrightarrow\quad
 \{E_\mu,E_\nu\}_\sharp,\ [D_\mu,D_\nu].
\]
Antikomutátor nese metriku, zatímco komutátor kovariantních derivací nese
křivost. V minimální lokální konstrukci není zapotřebí vkládací mapa,
průměrování přes kompaktní vlákno ani projekce pozorovatele.

\section{Obyčejná, parciální a kovariantní derivace}
Pro běžnou skalární funkci $f(x)$ derivace $df/dx$ měří, jak se
hodnota $f$ změní při změně $x$. Pro pole na časoprostoru používáme
parciální derivace,
\[
 \partial_\mu\Theta:=\frac{\partial\Theta}{\partial x^\mu},
 \qquad \mu=0,1,2,3.
\]
Parciální derivace porovnává složky zapsané v pevné bázi. To stačí,
lze-li všude použít tutéž bázi. Geometrie je jemnější, když se lokální
rámec při přechodu z bodu do bodu sám otáčí nebo Lorentzovsky transformuje.

Jednoduchou analogií je vektor zapsaný v rotující souřadnicové soustavě.
I fyzikálně konstantní vektor v ní má proměnné složky, protože se mění
báze. Kovariantní derivace tuto změnu báze koriguje:
\[
 D_\mu\Theta=\partial_\mu\Theta+\rho_*(\Omega_\mu)\Theta.
\]
Zde je $\Omega_\mu$ konexe a $\rho_*$ určuje, jak konexe působí v
reprezentaci pole $\Theta$. U biquaternionů jsou levá, pravá a oboustranná
akce skutečně odlišné a nesmějí se mlčky zaměňovat.

\section{Co je konexe?}
Konexe sama o sobě není silou. Je pravidlem pro porovnávání lokálních
rámců v sousedních bodech. Symbol $\Omega_\mu$ označuje interní neboli
biquaternionickou rámcovou konexi UBT.

Christoffelovy symboly
\[
 \Gamma^\rho{}_{\mu\nu},
\]
s ní úzce souvisejí, ale působí na souřadnicové indexy. Například
\[
 \nabla_\mu V^\rho=\partial_\mu V^\rho
 +\Gamma^\rho{}_{\mu\sigma}V^\sigma.
\]
Tedy:
\begin{itemize}
 \item $\Gamma^\rho{}_{\mu\nu}$ nám říká, jak se mění souřadnicová báze;
 \item $\omega_\mu{}^a{}_b$ a její spinový/biquaternionický obraz $\Omega_\mu$
       určují, jak se mění lokální Lorentzův rámec.
\end{itemize}
Jsou spojeny identitou kompatibility tetrád
\[
 \partial_\mu e_\nu{}^a-\Gamma^\rho{}_{\mu\nu}e_\rho{}^a
 +\omega_\mu{}^a{}_b e_\nu{}^b=0.
\]
Jde tedy o dva popisy téhož geometrického transportu, nikoli o tytéž
koeficienty. Zejména $\Gamma=\operatorname{Re}\Omega$ není platná
identifikace.

\section{Proč klasická konexe není další volné pole}
Pro nedegenerovanou tetrádu určují metrická kompatibilita a nulová torze
jedinou Lorentzovu konexi:
\[
 \boxed{
 \mathring\omega_\mu{}^a{}_b
 =e_b{}^\nu\left(\mathring\Gamma^\rho{}_{\mu\nu}(g)e_\rho{}^a
 -\partial_\mu e_\nu{}^a\right).}
\]
Kroužek označuje beztorzní Leviho--Civitovu konexi. V obvyklé beztorzní
větvi GR se tedy rámcová konexe vypočítá z tetrády, stejně jako se
Christoffelovy symboly počítají z metriky. Lokální Lorentzova transformace
změní číselné koeficienty tetrády i konexe, nikoli samotnou geometrii.

\subsection{Torze a kontorze}
Torzní dvoutvar je
\[
 T^a=de^a+\omega^a{}_b\wedge e^b.
\]
Lze zapsat obecnou metricky kompatibilní konexi
\[
 \boxed{\omega=\mathring\omega(e)+K(T),}
\]
kde $K$ je kontorze. Pro
$T^a=\tfrac12T^a{}_{bc}e^b\wedge e^c$ a výše uvedenou konvenci platí
\[
 \boxed{
 K_{abc}=\frac12\left(T_{cab}-T_{abc}-T_{bca}\right).}
\]
Tento vzorec má důležitý koncepční důsledek: jakmile jsou zadány tetráda
a torze, metricky kompatibilní konexe je jednoznačná. Otevřenou otázkou plné
UBT tedy není libovolný zbytek $\Omega_\mu$, ale dynamický zákon vybírající
$T$.

Pokud vakuové rovnice UBT implikují $T=0$, konexe se redukuje na
Leviho--Civitovu spinovou konexi. Pokud spin nebo jiný biquaternionický zdroj
vytvoří torzi, tatáž věta rekonstruuje konexi z této torze.

\section{Plochý prostor a explicitní UBT reprezentant}
V plochém Minkowského prostoročase lze v inerciálních kartézských
souřadnicích a konstantním lokálním rámci zvolit
\[
 \Gamma^\rho{}_{\mu\nu}=0,
 \qquad \Omega_\mu=0,
 \qquad D_\mu=\partial_\mu.
\]
Toto tvrzení lze posílit od pouhé volby kalibrace k explicitnímu
reprezentantu tvořenému jediným polem:
\[
 \boxed{
 \Theta_{\mathrm{SR}}(x)=\Theta_0+\sqrt{\mathcal N_0}
 \left(i x^0\mathbf1+x^k\mathbf e_k\right).}
\]
Pak
\[
 \mathcal N_0^{-1/2}\partial_0\Theta_{\mathrm{SR}}=i\mathbf1,
 \qquad
 \mathcal N_0^{-1/2}\partial_k\Theta_{\mathrm{SR}}=\mathbf e_k.
\]
Centrální antikomutátor těchto čtyř konstantních směrů dává
Minkowského metriku. Všechny druhé časoprostorové derivace
$\Theta_{\mathrm{SR}}$ mizí, takže pole v této ploché větvi řeší i
standardní bezzdrojový hlavní operátor druhého řádu s konstantními
koeficienty.

Obecněji řečeno, pro jakoukoli konstantní Lorentzovu tetrádu $\bar E_\mu$,
\[
 \Theta_{\mathrm{aff}}(x)=\Theta_0+
 \sqrt{\mathcal N_0}\,\bar E_\mu x^\mu
\]
generuje přesně tuto tetrádu. Tím se uzavírá problém reprezentace ve
speciální relativitě, nikoli však jedinečnost ani existence v zakřiveném
prostoru.

V polárních souřadnicích nebo ve zrychleném rámci mohou být $\Gamma$ a
$\Omega$ nenulové, i když je časoprostor plochý. Invariantním tvrzením o
plochosti je nulová křivost, nikoli nulové koeficienty konexe v každém rámci.

\section{Čtyři kovariantní gradienty jako tetráda}
Definujme
\[
 E_\mu:=\frac{1}{\sqrt{\mathcal N_0}}D_\mu\Theta.
\]
Čtyři objekty $E_0,E_1,E_2,E_3$ jsou lokálními měřítky UBT a určují
směr hodin. Nejde o čtyři dodatečná hmotová pole, ale o první kovariantní
derivace jediného fundamentálního pole $\Theta$.

Nechť $i$ je obvyklá komutující komplexní jednotka a
$\mathbf e_1,\mathbf e_2,\mathbf e_3$ jsou kvaternionové jednotky. V klasickém
Lorentzově sektoru pišme
\[
 E_\mu=i e_\mu{}^0\mathbf1+e_\mu{}^k\mathbf e_k,
 \qquad e_\mu{}^a\in\mathbb R.
\]
Kvaternionová konjugace mění znaménka tří kvaternionových jednotek, ale
komplexní skalární koeficient ponechává beze změny:
\[
 E_\mu^\sharp=i e_\mu{}^0\mathbf1-e_\mu{}^k\mathbf e_k.
\]

\section{Proč je antikomutátor metrikou}
Kvaternionové násobení obsahuje skalární i vektorový součin. V
symetrizovaném součinu se vektorové součiny zruší:
\[
 \boxed{
 \frac12\left(E_\mu^\sharp E_\nu+E_\nu^\sharp E_\mu\right)
 =g_{\mu\nu}\mathbf1.}
\]
Pravá strana obsahuje $\mathbf1$, jednotkový biquaternion. Celá levá strana
je již centrální: jde o obyčejné reálné číslo vynásobené jednotkou
algebry. Není třeba stopa, zobrazení na reálnou část, fázová projekce ani
průměrování přes vlákno.

Rozepsání koeficientů dává
\[
 g_{\mu\nu}
 =-e_\mu{}^0e_\nu{}^0+
  e_\mu{}^1e_\nu{}^1+
  e_\mu{}^2e_\nu{}^2+
  e_\mu{}^3e_\nu{}^3,
\]
což je přesně
\[
 g_{\mu\nu}=e_\mu{}^a e_\nu{}^b\eta_{ab},
 \qquad \eta_{ab}=\operatorname{diag}(-1,1,1,1).
\]
Minusové znaménko času pochází z $i^2=-1$.

\section{Co dělá druhá polovina součinu?}
Antisymetrická algebraická část je
\[
 \Sigma_{\mu\nu}
 =\frac12\left(E_\mu^\sharp E_\nu-E_\nu^\sharp E_\mu\right).
\]
Popisuje orientovanou rovinu neboli bivektor a je přirozeným nosičem lokální
informace o rotaci, boostu a spinu. Sama o sobě však není křivostí.

Tři výrazy musí být odděleny:
\begin{align*}
 \{E_\mu,E_\nu\}_\sharp
 &\longrightarrow g_{\mu\nu},\\
 [E_\mu,E_\nu]_\sharp
 &\longrightarrow \Sigma_{\mu\nu},\\
 [D_\mu,D_\nu]
 &\longrightarrow \mathcal R_{\mu\nu}.
\end{align*}

\section{Křivost z nekomutujících derivací}
Pro jednostrannou konexi má křivost známý tvar
\[
 \mathcal R_{\mu\nu}
 =\partial_\mu\Omega_\nu-\partial_\nu\Omega_\mu
  +[\Omega_\mu,\Omega_\nu].
\]
Poslední komutátor je nezbytný, protože biquaternionové násobení je
nekomutativní. Geometricky křivost měří, jak se rámec změní po paralelním
přenosu kolem malé uzavřené smyčky.

Rovnice $E_\mu=D_\mu\Theta/\sqrt{\mathcal N_0}$ vytváří další podmínku
integrability. Právě zde se stává rozhodující reprezentace konexe.

\section{Proč starý problém s hodností lokálně mizí}
Symetrická metrika $4\times4$ má deset nezávislých složek. Pohled pouze na
jednu hodnotu $\Theta\in\mathbb B\simeq\mathbb R^8$ svádí k zavádějícímu
srovnání $8<10$.

Metrika se nekonstruuje pouze z hodnoty $\Theta$, ale ze čtyř kovariantních
derivací $E_\mu$. V Lorentzově sektoru tvoří reálnou tetrádu
$4\times4$, $e_\mu{}^a$, se šestnácti složkami. Šest změn odpovídá lokálním
rotacím a Lorentzovým transformacím, které metriku nemění. Proto
\[
 16-6=10.
\]
Přesněji: každá malá symetrická variace metriky $h_{\mu\nu}$ vznikne volbou
\[
 \delta e_\mu{}^a=\frac12 h_{\mu\rho}e^{\rho a}.
\]
Zobrazení tetrády na metriku má proto u každé nedegenerované tetrády
hodnost deset. Tím je vyřešena lokální kinematická otázka hodnosti. Zatím
to nedokazuje, že každou zakřivenou tetrádu generuje pole UBT splňující
pohybové rovnice.

\subsection{Správný výpočet může odpovědět na špatně položenou architektonickou otázku}
Dřívější obstrukce hodnosti pro jediný řez byla matematicky správná v
rámci indukované metriky, kde se měly Einsteinovy rovnice obnovit pouze z
normálových variací jednoho řezu podobného vnoření. Nepředstavovala však
vylučovací větu pro kovariantní tetrádu
$E_\mu=D_\mu\Theta/\sqrt{\mathcal N_0}$. Konstrukce s kompaktním vláknem
opravila dřívější formulaci přidáním profilových směrů, ale tetrádová
formulace ukazuje, že původní architektura tuto opravu nevyžadovala.

Toto rozlišení motivuje obecné pravidlo výzkumu:

\begin{quote}
Před přidáním polí, módů, rozměrů, projekcí, průměrů nebo pomocných
struktur k odstranění obstrukce je nejprve nutné ověřit, zda obstrukce není
artefaktem formulace, v níž byla odvozena.
\end{quote}

Vláknová cesta zůstává matematicky konzistentní ve svém vlastním rámci.
Jejím problémem je slabý výběr a redundance reprezentantů, nikoli algebraický
rozpor.

\section{Problém integrability}
Jakmile je klasická konexe rekonstruována z tetrády a torze, má definiční
rovnice schematický tvar
\[
 E_\mu=\mathcal N_0^{-1/2}
 \left(\partial_\mu\Theta+\rho_*(\Omega_\mu[E,T])\Theta\right).
\]
Neznámé $E$ se vyskytuje na obou stranách. Jde o implicitní nelineární
parciální diferenciální rovnici prvního řádu nebo o systém pevného bodu.
Kinematická eliminace $\Omega$ sama o sobě nedokazuje existenci ani
jednoznačnost páru $(\Theta,E)$.

\subsection{Jednostranná překážka vedoucí k plochosti}
Předpokládejme, že konexe působí pouze zleva,
\[
 D^L_\mu\Theta=\partial_\mu\Theta+A_\mu\Theta,
\]
takže
\[
 [D^L_\mu,D^L_\nu]\Theta=F^L_{\mu\nu}\Theta.
\]
V beztorzní větvi kompatibilní tetrády dává antisymetrizace kovariantního
Hessiánu
\[
 F^L_{\mu\nu}\Theta=0.
\]
Pokud je $\Theta$ invertibilní jako biquaternion, pak
\[
 F^L_{\mu\nu}=0.
\]
Naivní jednostranná regulární reprezentace tedy nemůže obecně podporovat
zakřivenou GR a současně zachovat invertibilní $\Theta$ i beztorzní
kompatibilitu tetrády. Konfigurace s neinvertibilním stabilizátorem se mohou
tomuto výsledku vyhnout, nelze však předpokládat, že reprezentují libovolnou
geometrii.

\subsection{Přirozená oboustranná bikvaternionová derivace}
Biquaterniony přirozeně připouštějí levé i pravé násobení. Uvažujme
\[
 \boxed{
 D_\mu\Theta=\partial_\mu\Theta+A_\mu\Theta-\Theta B_\mu.}
\]
Přímý výpočet dává
\[
 \boxed{
 [D_\mu,D_\nu]\Theta
 =F^A_{\mu\nu}\Theta-\Theta F^B_{\mu\nu}.}
\]
Beztorzní kompatibilita proto vyžaduje
\[
 F^A_{\mu\nu}\Theta=\Theta F^B_{\mu\nu}.
\]
Pro invertibilní $\Theta$,
\[
 \boxed{
 F^A_{\mu\nu}=\Theta F^B_{\mu\nu}\Theta^{-1}.}
\]
Na rozdíl od jednostranného případu nemusí mizet ani jedna křivost. Pole
$\Theta$ proplétá levou a pravou reprezentaci křivosti. Jde o skutečné
strukturální zúžení GAP-10I: oboustranná neboli bimodulová derivace je
minimální cesta přirozená dané algebře, která se vyhýbá obecné obstrukci
plochosti.

\subsection{Redukce bez nového pole, beztorzní vylučovací výsledek a lokální pravá inverze s torzí}
Na Lorentzově řezu $W_L=\{X\mid X^\ddagger=-X\}$ redukuje zachování řezu
a jeho kvadratické formy čistý pár, až na společnou centrální jednoformu,
která se zruší, na
\[
 \boxed{A_\mu=\Omega_\mu,\qquad B_\mu=-\Omega_\mu^\ddagger.}
\]
$A_\mu$ a $B_\mu$ tedy v této větvi nejsou nezávislá gravitační pole.
Jde o dvě modulové akce jediné spinové konexe.

Při nulové torzi implikuje tatáž derivace kompatibilní s involucí spolu s
nedegenerovaným řešením $D_\mu\Theta=\sqrt{\mathcal N_0}E_\mu$ vztah
\[
 \boxed{\mathring\nabla_\mu V^\nu=\delta_\mu{}^\nu,}
 \qquad \mathcal L_Vg=2g.
\]
Neplochý Schwarzschildův vakuový exteriér s nenulovou hmotností takovou
vlastní homotetii nemá. Čistý pár je tedy kinematicky uzavřen a jeho větev
$K=0$ je uzavřena jako beztorzní vylučovací výsledek.

Na tomto upřesnění záleží. Na dostatečně malé nenulové Gaussově oblasti
zvolme $\rho\ne0$ tak, aby
$g^{\mu\nu}\partial_\mu\rho\partial_\nu\rho=\varepsilon=\pm1$, a položme
\[
 V^\mu=\varepsilon\rho\nabla^\mu\rho,
 \qquad
 W_{\mu\nu}=g_{\mu\nu}-\mathring\nabla_\mu V_\nu,
\]
\[
 \boxed{K_{\nu\mu\rho}
 =\frac{W_{\mu\nu}V_\rho-V_\nu W_{\mu\rho}}{V^2}.}
\]
Pak $K_{\nu\mu\rho}=-K_{\rho\mu\nu}$ a
$\nabla^{(\mathring\Gamma+K)}_\mu V^\nu=\delta_\mu{}^\nu$.  Proto
\[
 \Theta=\sqrt{\mathcal N_0}V^a\mathbf u_a
\]
splňuje $D_\mu\Theta=\sqrt{\mathcal N_0}E_\mu$. Každá hladká tetráda má tedy
takového lokálního reprezentanta s kompozitní metricky kompatibilní kontorzí
a bez nezávislých polí $A_\mu,B_\mu$.

Relativní člen $P_\mu X-XQ_\mu$ již není nutný pro lokální kinematickou
existenci. Zůstává možným beztorzním doplněním, a pokud se použije, musí
být kompozitní nebo pomocný, nikoli novým propagujícím polem. Společná
centrální konstanta beztorzní obstrukci neodstraní, protože se vyruší.

\section{Podmíněná dynamická dílčí uzavření}
\subsection{Algebraická Cartanova rovnice torze}
V minimální Hilbertově--Palatiniho větvi prvního řádu dává nezávislá
variace Lorentzovy konexe
\[
 \epsilon_{abcd}T^a\wedge e^b=\kappa\tau_{cd}.
\]
Pro nedegenerovanou tetrádu jde o bodovou lineární mapu $24\times24$
hodnosti 24. Nulový spinový proud tedy dává $T=0$, zatímco zadaný spinový
proud určuje jednoznačnou torzi a kontorzi. Tím se uzavírá algebra torze
uvnitř minimální větve, nikoli odvození této větve z fundamentální UBT.

\subsection{Propagace Lorentzova řezu jako množiny pevných bodů}
Antilineární involuce
\[
 \mathcal JX=-\overline{X^\sharp}
\]
má fyzický Lorentzův modul jako množinu pevných bodů. Je-li korektně
formulovaný jednoznačný Cauchyho vývoj $\mathcal J$-ekvivariantní a jsou-li
jeho data a zdroje pevné vzhledem k $\mathcal J$, z jednoznačnosti plyne, že
řešení zůstává na Lorentzově řezu. Konečnou akci UBT je ještě nutné
porovnat s těmito hypotézami.

\subsection{Stabilita metriky v imaginárním čase}
Vertikální tok
\[
 \partial_\psi e_\mu{}^a=\Lambda_\psi{}^a{}_b e_\mu{}^b,
 \qquad \Lambda_{\psi ab}=-\Lambda_{\psi ba},
\]
je tečný k lokální Lorentzově kalibrační orbitě a identicky splňuje
$\partial_\psi g_{\mu\nu}=0$. Druhým postačujícím mechanismem je jednoznačný
vývoj invariantní vůči posunům v $\psi$ s počátečními daty nezávislými na
$\psi$. Výběr fyzického mechanismu a vyloučení nestabilních nekalibračních
módů zůstává dynamickým úkolem.

\subsection{Rozšířená holonomie pro předepsané zakřivené koeficienty}
Pro předepsané $E_\mu,A_\mu,B_\mu$ se nehomogenní systém
\[
 \partial_\mu\Theta+A_\mu\Theta-\Theta B_\mu
 =\sqrt{\mathcal N_0}E_\mu
\]
mění na homogenní paralelní transport pro $Y=(\Theta,1)^T$ v rozšířeném
svazku. Jednohodnotové řešení s počáteční hodnotou $Y_0$ existuje právě tehdy,
když rozšířená holonomie ponechává $Y_0$ pevné; na jednoduše souvislé
oblasti stačí nulová rozšířená křivost. Tím se uzavírá integrabilita pro
předepsané koeficienty, nikoli jejich samokonzistentní výběr v UBT.

\subsection{Podmíněný Einsteinův--Lambda koncový bod}
Minimální Palatiniho větev s nulovým spinovým proudem dává
\[
 G_{\mu\nu}+\Lambda g_{\mu\nu}=\kappa T_{\mu\nu}.
\]
Za obvyklých čtyřrozměrných Lovelockových předpokladů má každá lokální
přirozená metrická rovnice druhého řádu s nulovou divergencí a bez dalších
lehkých geometrických polí gravitační část
$aG_{\mu\nu}+bg_{\mu\nu}$. Podmíněný nízkoenergetický koncový bod je proto
jednoznačný. UBT však musí tyto předpoklady, koeficienty a hmotové členy
teprve odvodit ze své kanonické akce.

\section{Implicitní versus transcendentní rovnice}
Slova \emph{implicitní} a \emph{transcendentní} popisují různé
vlastnosti.

Rovnice je implicitní, když se neznámá vyskytuje uvnitř zobrazení, které ji
má určit:
\[
 E=\mathcal F(E,\Theta).
\]
Rovnice je transcendentní, když obsahuje nealgebraické funkce, například
exponenciálu, logaritmus, goniometrickou funkci nebo Jacobiho theta funkci:
\[
 E=e^{-E}.
\]
Tetrádová rovnice UBT je již implicitní a diferenciální. Je-li přípustné
pole $\Theta(q,\tau)$ omezeno na třídu Jacobiho theta funkcí nebo jiných
transcendentních funkcí, může být úplný systém implicitní i transcendentní.
Ve formálním textu je třeba tyto vlastnosti uvádět odděleně, i když obě
intuitivně působí dojmem, že je neznámá ``zapletena sama do sebe''.

\section{Derivace versus variace}
Derivace, například $D_\mu\Theta$, popisuje, jak se konfigurace jednoho pole
mění v časoprostoru. Variace porovnává blízké možné konfigurace:
\[
 \Theta_\varepsilon=\Theta+\varepsilon\,\delta\Theta.
\]
Variace je
\[
 \delta\Theta=
 \left.\frac{d\Theta_\varepsilon}{d\varepsilon}\right|_{0}.
\]
Variace jsou zapotřebí při odvozování rovnic pole z akce. Nenahrazují
časoprostorové derivace. Závisí-li konexe na poli,
\[
 \delta(D_\mu\Theta)
 =D_\mu(\delta\Theta)+\rho_*(\delta\Omega_\mu)\Theta
\]
a je nutné zahrnout i indukovanou variaci konexe.

\section{Co je uzavřeno a co zůstává otevřené}
Následující výsledky jsou uzavřeny na uvedených úrovních:
\begin{enumerate}
 \item centrální antikomutátorová metrika a zobrazení tetrády na metriku
       hodnosti deset;
 \item GAP-10$\Omega$-KIN/GR: rekonstrukce kompatibilní konexe a její
       beztorzní Leviho--Civitova větev;
 \item GAP-10T-PALATINI (podmíněné): minimální Cartanova torzní rovnice je
       algebraická a invertibilní;
 \item GAP-10L-CONN a GAP-10L-SYM (podmíněné): kompatibilní transport a
       jednoznačný ekvivariantní vývoj zachovávají Lorentzův řez;
 \item GAP-10I-SR a GAP-10I-PRESCRIBED: afinní plochý zástupce a
       přesné kritérium rozšířené holonomie pro předepsané zakřivené koeficienty;
 \item GAP-10I-1S: naivní jednostranná invertibilní zakřivená cesta je
       podmíněný vylučovací výsledek;
 \item GAP-10I-PAIR-KIN: čistý kompatibilní pár se redukuje na jedinou
       spinovou konexi;
 \item GAP-10I-PAIR-GR: čistý pár s $K=0$ vede k vylučovacímu výsledku se souběžným
       vektorem pro neplochý Schwarzschildův vakuový exteriér s nenulovou
       hmotností;
 \item GAP-10I-TORSION-LOCAL: každá hladká tetráda má lokální jedno-$\Theta$
       reprezentanta s explicitní kompozitní metricky kompatibilní kontorzí;
 \item GAP-10D-PALATINI/UNIQUENESS (podmíněné): minimální první řád
       větev a Lovelockovy hypotézy vybírají Einsteinův--$\Lambda$ koncový bod;
 \item GAP-10$\psi$-KIN/SYM: Lorentzův kalibrační tok nebo translační symetrie
       chrání klasickou metriku podle uvedených hypotéz.
\end{enumerate}

Nerozdělené mezery plné teorie jsou zúžené, nikoli zcela uzavřené:
\begin{enumerate}
 \item GAP-10T-SPIN je podmíněně uzavřen pro přímou Palatiniho variaci na
       pevném pozadí, zatímco GAP-10T-FLAT-NOGO a
       GAP-10T-PAIRING-NOGO jsou uzavřeny jako vylučovací výsledky;
 \item GAP-10T-DYN: vypočítejte plnou složenou variaci $\Theta$ a odvoďte
       vybranou větev prvního řádu, kanonické zrušení torze nebo
       translační/relativní doplnění a normalizaci z kanonické UBT;
 \item GAP-10I-CURVED: místní kinematika je uzavřena; odvodit kanonický
       výběr na úrovni akce a fyzickou přípustnost kompozitní torzní větve
       nebo volitelné beztorzní relativní bimodulové větve a stanovit
       regularitu a globální pokračování
       $(\Theta,E,A,B,T)$;
 \item GAP-10L-DYN: dokázat ekvivarianci a korektně formulovanou jednoznačnost
       konečné úplné dynamiky;
 \item GAP-10D: odvodit nízkoenergetické Palatiniho/Lovelockovy předpoklady,
       vazby a hmotovou akci z fundamentální teorie;
 \item GAP-10$\psi$: vyberte mechanismus stability a vylučte fyzické
       nestabilní nekalibrační módy.
\end{enumerate}
GAP-B-MASTER a GAP-U2$\Theta$ zůstávají otevřené. Dřívější uzávěra s
kompaktním vláknem $\psi$ zůstává matematicky konzistentním historickým
kandidátním doplněním, ale její slabý výběr a redundance reprezentantů ji
udržují mimo minimální kanonickou cestu.
